\documentclass[11pt]{article}
\usepackage[margin=1in]{geometry}
\usepackage{amsmath,amssymb,amsthm,array}
\usepackage[colorlinks=true,linkcolor=blue,citecolor=blue,urlcolor=blue]{hyperref}
\newtheorem{theorem}{Theorem}[section]
\newtheorem{lemma}[theorem]{Lemma}
\newtheorem{proposition}[theorem]{Proposition}
\newtheorem{corollary}[theorem]{Corollary}
\newtheorem{remark}[theorem]{Remark}
\numberwithin{equation}{section}
\newcommand{\R}{\mathbb R}
\newcommand{\C}{\mathbb C}
\newcommand{\rank}{\operatorname{rank}}
\newcommand{\diag}{\operatorname{diag}}
\newcommand{\sgn}{\operatorname{sgn}}
\newcommand{\tr}{\operatorname{tr}}
\newcommand{\op}{\mathrm{op}}
\newcommand{\Frob}{\mathrm F}
\newcommand{\nn}{\mathcal N}
\title{Sharp-rank algorithms for low-rank updates\texorpdfstring{\\}{ }of fractional matrix functions}
\author{Working research manuscript}
\date{November 2026}
\begin{document}
\maketitle

\begin{abstract}
We give algorithms for low-rank corrections to fractional matrix
functions from thin raw factors. The error is
$\varepsilon\|E\|_{\op}^{\alpha}$ for the actual perturbation $E$,
including arbitrarily cancelling supplied factors. Conventional arithmetic
cost is $\widetilde O_\alpha(Nmr+m^3)$, with no condition number,
spectral gap, or supplied base matrix function. For fixed perturbation
rank $r$ and $\varepsilon\downarrow0$, the output width has the sharp
leading coefficient in the uniform worst-case rank bound. At
$\alpha=1/2$ the coefficients of $r\log^2(1/\varepsilon)$ are
$1/\pi^2$ for positive-endpoint roots, $2/\pi^2$ for a complex branch,
and $4/\pi^2$ for real whole-line maps and rectangular absolute/polar
fractional maps. The square-root algorithms use only field operations,
comparisons, scalar square roots, and discrete randomness. We explicitly
realize the positive Gaussian quadrature without exponential or error
function oracles. For rectangular right absolute and polar corrections,
forming every output entry is included in the same asymptotic cost.
\end{abstract}

\noindent\textit{Review and priority status.}
The mathematical notes and complete typeset source passed separate
independent internal audits, including an actual-copy review. There is
no claim of external peer review, formal verification, finite-word
stability, or certified novelty. Rational matrix-function updates,
root-exponential compression, and randomized preconditioning have
substantial antecedents. The precise synthesis and sharp output
guarantees, rather than those individual tools, are the subject here.

\section{Input, guarantees, and sharp leading widths}

All input matrices are real, and $\|\cdot\|$ means the operator norm.
Fix $0<\alpha<1$ and $\beta=1-\alpha$. Costs count exact real field
operations, comparisons, positive scalar square roots, and discrete
randomness. Complex shifted solves use pairs of real coordinates.
The notation $\widetilde O$ suppresses only logarithms of dimensions,
$1/\varepsilon$ and $1/\delta$, where $\delta$ is the failure probability.
Constants may depend on the fixed scalar function and $\alpha$.

The principal case $\alpha=1/2$, and more generally dyadic $\alpha$,
uses the preceding scalar operations alone. For arbitrary fixed
$\alpha$, input-dependent scalar powers such as $a^\alpha$ must be
added to the model or their evaluation cost charged. Our statements
for such $\alpha$ include this qualification. Fixed scalar constants
can be approximated by ordinary convergent series at polynomial cost
in $\log(1/\varepsilon)$.

\begin{theorem}[Positive endpoints]\label{thm:PSD}
Let $X\in\R^{N\times m}$, $U\in\R^{N\times r}$,
$1\le r\le m\le N$, and let $K\in\R^{r\times r}$ be symmetric.
Set $A=XX^T$, $E=UKU^T$, and promise $B=A+E\succeq0$.
For $\varepsilon,\delta\in(0,1)$ an algorithm returns a real signed
factor $F=Z\Sigma Z^T$, where $\Sigma$ is diagonal with entries
$\pm1$, satisfying
\begin{equation}\label{eq:errorPSD}
 \Pr\{\|F-(B^\alpha-A^\alpha)\|
             \le\varepsilon\|E\|^\alpha\}\ge1-\delta
\end{equation}
in conventional arithmetic $\widetilde O_\alpha(Nmr+m^3)$.
Its width is $O_\alpha(r\log^2(2r/\varepsilon))$, and for each fixed
$r$, uniformly over dimensions and input spectra, the width is at most
\begin{equation}\label{eq:PSDwidth}
 \left[\frac{r}{4\pi^2\alpha\beta}+o_r(1)\right]
                          \log^2(1/\varepsilon).
\end{equation}
If $E\succeq0$, $\Sigma$ may be taken positive; if $E\preceq0$,
negative. Singular endpoints and critical downdates are included.
Dense expansion adds $O(N^2\,\mathrm{width})$.
\end{theorem}

For $X\in\R^{N\times m}$, $U\in\R^{N\times r}$,
$1\le r\le m\le N$, and symmetric $K\in\R^{r\times r}$,
write a general signed raw base as
\begin{equation}\label{eq:rawS}
 S=XJX^T,\qquad J=J^T\in\R^{m\times m},
 \qquad T=S+UKU^T.
\end{equation}
Neither $J$ nor $K$ must be nonsingular or an involution. Define
\[
 f_e(x)=|x|^\alpha,\quad
 f_o(x)=\sgn(x)|x|^\alpha,\quad
 g_\pm(x)=(\mp ix)^\alpha,
\]
using the principal scalar branch and value zero at zero.

\begin{theorem}[Whole-line functions]\label{thm:whole}
For $\varepsilon,\delta\in(0,1)$ and the input~\eqref{eq:rawS},
with no positivity promise, the same
conventional cost returns an approximation to $f(T)-f(S)$ with error
$\varepsilon\|E\|^\alpha$ and probability at least $1-\delta$.
For $f=f_e,f_o$, or any fixed nonzero real homogeneous function
$a_0x_+^\alpha+b_0(-x)_+^\alpha$, the output is a real signed factor
with leading width
\begin{equation}\label{eq:wholewidth}
 \left[\frac{r}{\pi^2\alpha\beta}+o_r(1)\right]
                              \log^2(1/\varepsilon).
\end{equation}
For $g_\pm$ a complex factor $PQ^T$ has leading width
\begin{equation}\label{eq:branchwidth}
 \left[\frac{r}{2\pi^2\alpha\beta}+o_r(1)\right]
                              \log^2(1/\varepsilon).
\end{equation}
All widths are $O_\alpha(r\log^2(2r/\varepsilon))$ without taking a
limit. The scalar-operation qualification above applies throughout.
\end{theorem}

\begin{theorem}[Rectangular changes]\label{thm:rect}
Let $M\in\R^{N\times m}$, $N\ge m$, and
$\Delta=UV^T$ with $U\in\R^{N\times r}$, $V\in\R^{m\times r}$,
$1\le r\le m$. For $\varepsilon,\delta\in(0,1)$ put
$M'=M+\Delta$. The differences of
\[
 |M|^\alpha=(M^TM)^{\alpha/2},\qquad
 |M^T|^\alpha=(MM^T)^{\alpha/2},\qquad
 \mathcal P_\alpha(M)=M|M|^{\alpha-1}
\]
have low-rank approximations with error $\varepsilon\|\Delta\|^\alpha$,
the probability/cost of Theorem~\ref{thm:whole}, and leading width
\eqref{eq:wholewidth}. The polar map is defined by singular values,
with zero on the zero singular space. Forming every entry of the right
absolute-power or polar correction is included in
$\widetilde O_\alpha(Nmr+m^3)$. Left absolute-power expansion costs
an additional $O(N^2\,\mathrm{width})$.
\end{theorem}

The coefficients for $\alpha=1/2$, divided by $r$, are
\[
\begin{array}{c|c}
 \text{map or endpoint class}&\text{coefficient of }\log^2(1/\varepsilon)\\ \hline
 \text{PSD power}&1/\pi^2\\
 \text{one complex branch}&2/\pi^2\\
 \text{real whole-line map}&4/\pi^2\\
 \text{rectangular absolute/polar fractional map}&4/\pi^2
\end{array}
\]
The companion compression theorem~\cite{Sharp,Extensions} gives matching
worst-case lower coefficients for these classes, for each fixed $r$.
Thus ``sharp'' refers to the leading \emph{output rank}, over unrestricted
finite dimensions. It asserts neither optimal runtime nor the exact
minimum rank for a fixed input or fixed positive $\varepsilon$.

\section{An intrinsic singular perturbation core}

Compute rank-revealing thin QR $U=QR$, where $Q^TQ=I_k$ and $k\le r$.
Set
\begin{equation}\label{eq:intrinsic}
 C=RKR^T,\qquad E=QCQ^T,\qquad a=\|C\|_{\Frob}.
\end{equation}
If $a=0$, return zero. Otherwise
$\|E\|=\|C\|\le a\le\sqrt r\|E\|$.
This uses $O(Nr^2+r^3)$ work. It does not introduce an eigennorm oracle,
and it removes arbitrarily large cancellation in the supplied factors.

For a resolvent $R_z=(S-zI)^{-1}$ put
$R'_z=(T-zI)^{-1}$, $V_z=R_zQ$, $M_z=Q^TR_zQ$.
The same notation is used at $z=-t$ for positive endpoints.

\begin{lemma}[Core and perturbation bounds]\label{lem:core}
At $z=it$ for Hermitian endpoints, or $z=-t$ for PSD endpoints, $t>0$,
\begin{align}
 H_z&=(I+M_zC)^{-1}=I-Q^TR'_zQC,&\|H_z\|&\le\mu:=1+a/t,
                                                        \label{eq:coreH}\\
 K_z&=CH_z=(I+CM_z)^{-1}C=K_z^T,&
 D_z:=R_z-R'_z&=V_zK_zV_z^T.             \label{eq:coreD}
\end{align}
Suppose a common symmetric linear inverse approximation obeys
$\|\widehat R_z-R_z\|\le\nu/t$, $\nu\le1$.
Define $\widehat V,\widehat M,\widehat H,\widehat K$ by the same formulas.
If $\nu(a/t)(1+a/t)\le1/2$, then $\widehat H$ exists,
$\widehat M$ and $\widehat K$ are symmetric, and
\begin{align}
 \|\widehat K\|&\le2a\mu,&
 \|\widehat K-K\|&\le2\nu a^2\mu^2/t,\\
 \|\widehat V\widehat K\widehat V^T-VKV^T\|
       &\le6\nu a\mu/t^2+2\nu a^2\mu^2/t^3.   \label{eq:nodeError}
\end{align}
No inverse of $C$ is used; each computed node has rank at most
$\rank C=\rank E$.
\end{lemma}
\begin{proof}
The determinant lemma and endpoint resolvent invertibility give
$I+MC$ invertible. Woodbury gives both identities
\eqref{eq:coreH}--\eqref{eq:coreD}; the bound follows from
$\|R'_z\|\le1/t$ and $\|C\|\le a$. The equality
$C(I+MC)^{-1}=(I+CM)^{-1}C$ proves transpose symmetry, including
singular $C$. The resolvents and $M,K$ are complex symmetric in the
complex case, so the transposes rather than adjoints are intentional;
$H$ need not be symmetric.

Now $\|\widehat M-M\|\le\nu/t$. Inverse perturbation gives
$\|\widehat H\|\le2\mu$ and
$\|\widehat H-H\|\le2\nu a\mu^2/t$.
Use $\|V\|\le1/t$, $\|\widehat V\|\le2/t$ and
$\|\widehat V-V\|\le\nu/t$. Expanding the two $V$ factors and then
the core gives coefficients $4+2$ and $2$ in~\eqref{eq:nodeError}.
The rank assertion follows because $\widehat H$ is invertible.
\end{proof}

\section{Raw inverse actions with one sketch}

A constant-distortion OSE forms $G_0=L^TL$ satisfying, with probability
at least $1-\delta$,
\begin{equation}\label{eq:OSE}
 (1-\eta_0)G_0\preceq G=X^TX\preceq(1+\eta_0)G_0,
 \qquad\eta_0=1/4.
\end{equation}
The support factor $L$ is full row rank. Standard sparse embeddings
give $\widetilde O(Nm+m^3)$ conventional setup. Explicitly,
\cite[Theorem 9]{NN} has sparsity
$\Theta(\log^3(d/\delta)/\epsilon_0)$ and
$O(d\log^8(d/\delta)/\epsilon_0^2)$ rows for a $d$-dimensional
space, with success $1-\delta$. Use fixed sufficiently small
$\epsilon_0$ to obtain~\eqref{eq:OSE}; rank deficiency causes no issue.
Only the sketched Gram is formed. The same event controls every later
node, RHS and iteration.

\subsection*{Positive base}
For $A=XX^T$, set $P=G_0+tI$ and
$F_t=P^{-1/2}(G-G_0)P^{-1/2}$, so $\|F_t\|\le1/4$.
For an even integer $\ell$ define
\[
 Q_\ell=P^{-1/2}\sum_{j=0}^{\ell-1}(-F_t)^jP^{-1/2},\qquad
 \widehat R_{-t}=t^{-1}(I-XQ_\ell X^T).
\]
The inverse square roots are analysis only: Cholesky of $P$ and ordinary
preconditioned Richardson implement the polynomial. Every true
$G$ action is $X^T(XW)$ on a thin block. Since
\[
 0\preceq(G+tI)^{-1}-Q_\ell
          \preceq4^{-\ell}(G+tI)^{-1},
\]
we have the exact common-operator bound
\begin{equation}\label{eq:positiveR}
 0\preceq\widehat R_{-t}-R_{-t}\preceq4^{-\ell}I/t,
 \qquad \widehat R_{-t}\succeq0.
\end{equation}
The polynomial and its physical output are symmetric.

\subsection*{Signed base and imaginary poles}
For $S=XJX^T$ set $K_0=LJL^T$. At $\zeta=it$ put $z=1/\zeta$ and
\[
 A_z=I-zJG,\quad A_{0,z}=I-zJG_0,\quad b=JX^Th.
\]
Starting at $w_0=0$, perform
\begin{equation}\label{eq:Richardson}
 w_{j+1}=w_j+A_{0,z}^{-1}(b-A_zw_j),\qquad
 \widehat R_\zeta h=-\zeta^{-1}(h+zXw_\ell).
\end{equation}
The actual coarse inverse is
\[
 A_{0,z}^{-1}=I+zJL^T(I-zK_0)^{-1}L.
\]
It uses an $m$-dimensional complex solve; $I-zK_0$ is nonsingular
because $K_0$ is real symmetric and $z$ purely imaginary. Per-node
setup is $O(m^3)$ conventional arithmetic.

To prove convergence, write $G-G_0=L^TDL$ analytically with
$\|D\|\le\eta_0$. For the exact solution $w_*$ the recurrence is
\[
 L(w_{j+1}-w_*)=zK_0(I-zK_0)^{-1}DL(w_j-w_*).
\]
The first factor has norm at most one. Moreover
$\|[(I-zS)^{-1}-I]h\|\le\|h\|$. The Gram sandwich transfers
coarse seminorm error to physical error and gives
\begin{equation}\label{eq:imagR}
 \|\widehat R_{it}-R_{it}\|
       \le\sqrt{5/3}\,4^{-\ell}/t.
\end{equation}
This holds for the actual linear operator defined by the fixed iteration,
uniformly for every $h$; nullspace components of the latent iterate do
not matter.

Symmetry must also be checked. The matrix
\[
 B_{0,z}=A_{0,z}^{-1}J
       =J+zJL^T(I-zK_0)^{-1}LJ
\]
is complex symmetric. The finite inverse polynomial is
$\sum_{j=0}^{\ell-1}z^jB_{0,z}[(G-G_0)B_{0,z}]^j$.
Each alternating word begins and ends with $B_{0,z}$ and is transpose
symmetric. Thus $\widehat R_{it}$ is complex symmetric, and real
coefficients give $\widehat R_{-it}=\overline{\widehat R_{it}}$.
A generic RHS-dependent stopping rule would not supply this property.

\section{Integral representations and uniform envelopes}

For PSD endpoints and $c_\alpha=\sin(\pi\alpha)/\pi$,
\begin{equation}\label{eq:PSDintegral}
 B^\alpha-A^\alpha
       =c_\alpha\int_0^\infty t^\alpha D_{-t}\,dt.
\end{equation}
Here $\|D_{-t}\|\le\min(1/t,a/t^2)$, since both endpoint resolvents
lie between zero and $I/t$. For Hermitian endpoints,
$\|D_{it}\|\le\min(2/t,a/t^2)$. Put
$\phi=\pi\alpha/2$, $c_e=2\sin\phi/\pi$, $c_o=2\cos\phi/\pi$.
The whole-line formulas are
\begin{align}
 f_e(T)-f_e(S)&=c_e\int_0^\infty t^\alpha\Im D_{it}\,dt,
                                                   \label{eq:evenIntegral}\\
 f_o(T)-f_o(S)&=-c_o\int_0^\infty t^\alpha\Re D_{it}\,dt,\\
 g_+(T)-g_+(S)&=i c_\alpha\int_0^\infty t^\alpha D_{-it}\,dt.
                                                   \label{eq:branchIntegral}
\end{align}
The conjugate formula gives $g_-$. These identities follow from scalar
Stieltjes formulas and functional calculus; the displayed norm bounds
justify subtracting and integrating, including at zero eigenvalues.

Set $t=a e^y$ and divide the logarithmic integrand by $a^\alpha$.
For positive endpoints it is analytic on $|\Im y|<\pi$, with
\begin{equation}\label{eq:envelope}
 \|F(x+i\theta)\|\le M_d e^{-g(x)},\qquad
 g(x)=\alpha(-x)_++\beta x_+,
 \quad M_d\le C_\alpha\cos(d/2)^{-2}.
\end{equation}
Indeed $|\lambda+t|\ge\cos(d/2)(\lambda+|t|)$ for $\lambda\ge0$.
For whole-line maps the strip is $|\Im y|<\pi/2$, with
$M_d\le C_\alpha\cos(d)^{-2}$. Both imaginary resolvents then have
distance at least $a e^x\cos d$ from the real spectrum. For the analytic
extension of~\eqref{eq:evenIntegral}, use
$(D_{it}-D_{-it})/(2i)$, rather than interpreting $\Im$ off the real
line; similarly use their half-sum for the odd map.

On the real line $\|F(x)\|+\|F'(x)\|\le C_\alpha$ uniformly in
spectra and dimensions. For example
$D_t'=-R_AD_t-D_tR_B$ at negative real shifts, and the analogous
imaginary derivative has the same norm bound
$\|D_t'\|\le2\|D_t\|/t$. Differentiating the prefactor proves the
claim. These envelopes involve neither $\|S\|/\|E\|$ nor conditioning.

\section{A positive quadrature with explicit scalar realization}

We use the precise Gaussian localization result from the companion
paper~\cite[Lemma 3.1]{Sharp}. If $F$ is analytic on a strip of
half-width $d_0$ and satisfies~\eqref{eq:envelope}, then for fixed
$0<d<d_0$ there is a positive quadrature of error $e^{-L}$ and
\begin{equation}\label{eq:fixedGaussian}
 q\le\frac{L^2}{4\pi d}\left(\frac1\alpha+\frac1\beta\right)
                   +O_{\alpha,d}(L^{7/4}).
\end{equation}
We next make its moving-strip version and scalar operations explicit.
This step is needed for an algorithm attaining the exact leading rank.

\begin{lemma}[Constructive moving strip]\label{lem:constructive}
For the normalized integrands above, let $d_0=\pi$ or $\pi/2$.
For every sufficiently large $L$ one can compute positive nodes $s_j$
and nonnegative coefficients $a_j$, in $O_\alpha(L^3)$ scalar
arithmetic/square-root operations, such that
\begin{equation}\label{eq:movingCount}
 q\le\frac{L^2}{4\pi d_0\alpha\beta}+O_\alpha(L^{7/4}),
 \quad |\log s_j|\le C_\alpha(L+1),\quad
 \sum_j a_j\le C_\alpha(L+1).
\end{equation}
Replacing the ideal logarithmic quadrature by the actual nodes and
coefficients has normalized error at most $e^{-L}$ after fixed
constant budget adjustments. There is no exponential, Gaussian integral,
or error function oracle. At general $\alpha$ only the separately
qualified scalar powers in evaluating $F$ remain.
\end{lemma}
\begin{proof}
Compute certified rational intervals for $\pi$ and $\log2$ of width
$e^{-C_\alpha L}$, using Machin's arctangent identity and the
$\operatorname{atanh}(1/3)$ series. Their geometric tails cost
$O_\alpha(L)$ operations. Let $p\le\pi$ be the lower approximation;
put $d=p-L^{-1/4}$ or $p/2-L^{-1/4}$, respectively. For large $L$,
$M_d=O_\alpha(\sqrt L)$ by~\eqref{eq:envelope}. Use
\[
 \ell=L^{3/4}+C_\alpha,\quad b=2\sqrt L+C_\alpha,\quad
 I_j=[j\ell,(j+1)\ell],\quad
 \varphi_j(z)=\pi^{-1/2}\int_{I_j}e^{-(z-u)^2}\,du.
\]
Retain intervals meeting
$[-(L+\ell)/\alpha,(L+\ell)/\beta]$. Write
$m_j=\min_{I_j}g$, $Q_j=\max(\ell,L-m_j+3\ell)$ and
$h_j=2pd/Q_j$. Retain lattice points $x=nh_j$ in $I_j+[-b,b]$.
The ideal positive coefficient is $h_j\varphi_j(x)$.

For clarity, the companion proof's errors remain uniform as $d$ moves.
The boundary integrals of $F\varphi_j$ are at most
$C_\alpha M_d\ell e^{-m_j}$. Omitted windows contribute
$O_\alpha(M_de^{-L-\ell})$; infinite trapezoid errors sum to
$O_\alpha(M_dL e^{-L-3\ell})$, since $2\pi d/h_j\ge Q_j$.
Gaussian lattice tails outside the enlarged intervals contribute
$O_\alpha(M_dL e^{-2L})$. Polynomial growth $M_d=O_\alpha(\sqrt L)$
is absorbed by these margins. Fixed buffers or a finite-small-$L$
fallback make the desired error budget valid. All constants here come
from the displayed sector and Gaussian bounds and are computable for
fixed $\alpha$. One may choose a fixed threshold $L_\alpha$ from those
bounds and replace smaller requested $L$ by $L_\alpha$; this does not
change the asymptotic coefficient or introduce input spectral data.

Counting these lattices uses
\[
 \int_\R(L-g(x))_+\,dx
       =\frac{L^2}{2}\left(\frac1\alpha+\frac1\beta\right).
\]
The Riemann-sum and $3\ell$ losses are $O_\alpha(L\ell)$, and window
enlargement costs $O_\alpha(L^2b/\ell)$. Thus
$q\le L^2/(4pd\alpha\beta)+O_\alpha(L^{7/4})$, proving
\eqref{eq:movingCount}. Also $h_j<1$ for large $L$, and
$h\sum_n e^{-(nh-u)^2}<4$ for $h\le1$. Integrating over the retained
intervals gives total coefficient mass $O_\alpha(L+1)$.

Here is a positive scalar implementation of each window. Clip
$I_j-x$ to $[-R_0,R_0]$, where $R_0=4\sqrt{L+1}$. Set
\[
 D=\lceil64(L+1)\rceil,\qquad
 P_D(w)=\sum_{h=0}^D\frac{(-1/2)^h w^{2h}}{h!},\qquad
 \widehat\varphi_j(x)=p^{-1/2}\int_{(I_j-x)\cap[-R_0,R_0]}P_D(w)^2\,dw.
\]
The integral is a polynomial primitive, evaluated at two clipped
endpoints. It is nonnegative exactly. On the clipped interval,
$z=w^2/2\le8(L+1)$, so
\[
 |P_D(w)-e^{-w^2/2}|\le\frac{z^{D+1}}{(D+1)!}
       \le(e/8)^{D+1}\le2^{-D}.
\]
Squaring changes the error by at most $3\cdot2^{-D}$; the clipped
Gaussian tail is $O(e^{-16(L+1)})$. The $p$ error is smaller after
choosing its fixed accuracy constant. Consequently every window has
absolute error $O(e^{-16(L+1)})$. Precomputing $P_D^2$ costs $O(L^2)$,
and each primitive evaluation costs $O(L)$, for $O_\alpha(L^3)$ total.
Absolute accuracy suffices even for very small window weights.

Finally $|x|\le B=O_\alpha(L+1)$. Choose a power of two $M\ge8(B+1)$
of order $B+1$; compute a degree-$O_\alpha(L)$ Taylor approximation
$P$ to $e^{x/M}$, and set $s=P^M$ by repeated squaring. Then $P>0$,
$s>0$, and $|\log s-x|\le e^{-C_\alpha L}$. Taylor remainder at
$|x/M|\le1/8$ proves this estimate; the logarithm is analysis only.
The operation count is $O_\alpha(L)$ per node and independent of the
magnitude of the input scale $a$. By the real-line derivative bound
and the total coefficient mass, these node changes, the window errors,
and fixed scalar-constant approximations contribute at most $e^{-L}$
after a harmless constant shift of $L$. This proves the lemma.
\end{proof}

\section{The actual node algorithms and their costs}

Choose $\rho=\varepsilon/[C_\alpha(r+1)]$. Repeated halving finds the
least integer $n$ with $2^{-n}\le\rho$ in $O(\log(1/\rho))$ operations.
Before defining $L$, compute the constant intervals to width
$2^{-C_\alpha(n+1)}$ using the geometric series above; $O_\alpha(n+1)$
terms suffice. This supplies all requested constant precision without
a circular dependence on $L$.
Using the certified upper approximation to $\log2$, take
$L=n(\log2)_+$ plus a fixed budget shift. Thus
\begin{equation}\label{eq:Lscale}
 L=\log(1/\varepsilon)+O_\alpha(\log(2r)),\qquad e^{-L}\le\rho.
\end{equation}
All lattice endpoints and integer ceilings in the preceding construction
are found by comparisons within polynomial ranges in $L$; a floor or
logarithm of an input magnitude is not a free primitive.

For each actual node $t_j=as_j$, use the appropriate common inverse
from Section 3, with one fixed even iteration count $\ell_{\rm inv}$.
Let $B=\max_j|\log s_j|+1=O_\alpha(L+1)$ analytically; a computable
bound $1+\max_j|x_j|$ plus a fixed rounding allowance suffices. Choose
the least even integer $\ell_{\rm inv}\ge C_\alpha(B+L+1)$ large enough that
$\nu\le e^{-C_\alpha(B+L+1)}$ in~\eqref{eq:positiveR} or
\eqref{eq:imagR}. No spectral ratio enters this choice.

For example, multiplying~\eqref{eq:nodeError} by $t^{\alpha+1}$ and
putting $s=t/a$ gives an error at most a fixed multiple of
\[
 \nu a^\alpha\{s^{\alpha-1}(1+1/s)
                         +s^{\alpha-2}(1+1/s)^2\}.
\]
For $e^{-B}\le s\le e^B$ this is at most
$C_\alpha\nu a^\alpha e^{4B}$. The coefficient mass is
$O_\alpha(L+1)$, so the preceding fixed iteration count makes the
total solve error at most $a^\alpha e^{-L}$. It also ensures the
invertibility condition in Lemma~\ref{lem:core} at every actual node.
Together with Lemma~\ref{lem:constructive} the final error is at most
$C_\alpha a^\alpha e^{-L}\le\varepsilon\|E\|^\alpha$, since
$a^\alpha\le r^{\alpha/2}\|E\|^\alpha$ and
$r^{\alpha/2}/(r+1)\le1$.

The physical PSD weight is
$w_j=a_j\widehat c_\alpha t_j^{\alpha+1}$, with
$\widehat c_\alpha>0$ a scalar approximation to $c_\alpha$ of the
budgeted relative accuracy. Whole-line nodes use
$a_jt_j^{\alpha+1}$ multiplying the real or imaginary cores in
\eqref{eq:evenIntegral}--\eqref{eq:branchIntegral}, with their fixed
scalar coefficients approximated to the same accuracy. The needed
scalar powers use the operation model stated in Section 1.

\subsection*{Real and complex node factors}
For PSD endpoints $\widehat K$ is real symmetric with rank at most $r$.
Exact congruence elimination gives
$\widehat K=L_j\Sigma_jL_j^T$ with at most $\rank C$ columns.
Always search all remaining diagonal entries for a nonzero pivot first;
if all vanish and an off-diagonal $a_{ij}\ne0$ remains, replace
$e_i$ by $e_i+e_j$, producing diagonal $2a_{ij}$. Stop at a zero
residual matrix. This is $O(r^3)$ arithmetic, with square roots for
absolute diagonal pivots, and works for singular $C$.
The node factor is $\sqrt{w_j}\widehat V_jL_j$, where
$w_j$ is the nonnegative actual scalar weight.

If $C\succeq0$, then $\widehat M\succeq0$ by
\eqref{eq:positiveR}, and
$\widehat K=C^{1/2}(I+C^{1/2}\widehat MC^{1/2})^{-1}C^{1/2}\succeq0$.
These square roots are analysis only; PSD Cholesky with zero pivots
implements the factor. For $C\preceq0$ and $B\succeq0$, the exact
downdate core inverse has norm at most $1+a/t$. The same perturbation
condition keeps its approximated core positive, giving
$\widehat K\preceq0$. This proves the sign assertion in
Theorem~\ref{thm:PSD}.

For an imaginary node write $\widehat V=P+iQ_0$ and
$\widehat K=A_0+iB_0$, with real symmetric $A_0,B_0$. Then
\begin{align*}
 \Re\widehat D&=[P,Q_0]
       \begin{pmatrix}A_0&-B_0\\-B_0&-A_0\end{pmatrix}[P,Q_0]^T,\\
 \Im\widehat D&=[P,Q_0]
       \begin{pmatrix}B_0&A_0\\A_0&-B_0\end{pmatrix}[P,Q_0]^T.
\end{align*}
Each real node has rank at most $2r$. For a real homogeneous map,
combine these two cores on the \emph{same} feature block before
factoring; this avoids an additional factor two. A single complex
branch uses $\widehat V\widehat K\widehat V^T$ directly, at width
at most $r$.

The node count~\eqref{eq:movingCount}, multiplied by these node ranks,
proves~\eqref{eq:PSDwidth}--\eqref{eq:branchwidth} using~\eqref{eq:Lscale}.
All factors can be concatenated; no unknown matrix function is evaluated.

\subsection*{Operation and failure accounting}
Before the nodes, intrinsic QR/core costs $O(Nr^2+r^3)$, and the
sketch/coarse setup costs $\widetilde O(Nm+m^3)$. Reusable $X^TQ$
costs $O(Nmr)$. Each of $O_\alpha(L^2)$ nodes costs
$O(m^3+\ell_{\rm inv}Nmr+Nr^2+r^3)$; the scalar preparation costs
$O_\alpha(L^3)$. Thus, suppressing dimension/confidence logarithms only,
an explicit accuracy-dependent bound is
\[
 \widetilde O(Nm+m^3)+O_\alpha(NmrL^3+m^3L^2+L^3).
\]
This is $\widetilde O_\alpha(Nmr+m^3)$ with conventional multiplication.
Physical factor storage is $O_\alpha(NrL^2)$ in addition to input and
sketched-core storage. Output expansion has the costs already stated.

The only random failure allowance is the one OSE event. Every loop
count is fixed in advance. A singular small inverse or a failed required
PSD test returns zero; neither can occur on that event. Congruence zero
residuals represent rank loss and are not failure. Coarse imaginary
solves are always nonsingular, and positive $G_0+tI$ is always SPD.
Hence the operation/width cap holds on every random outcome. Verifying
the promised PSD endpoints is not an additional task of the algorithm.

\section{Rectangular dilation and a rank-preserving parity}

For $M\in\R^{N\times m}$ define
\[
 \mathcal D(M)=\begin{pmatrix}0&M\\M^T&0\end{pmatrix}
       =YJ_sY^T,\qquad
 Y=\diag(M,I_m),\qquad J_s=\begin{pmatrix}0&I_m\\I_m&0\end{pmatrix}.
\]
The latent width is $2m$, ambient size $N+m$, and raw input cost
$O(Nm+m^2)$. For $\Delta=UV^T$ its dilation is
$\diag(U,V)J_{s,r}\diag(U,V)^T$, of rank at most $2r$ and norm
exactly $\|\Delta\|$. Functional calculus gives
\begin{align*}
 f_e(\mathcal D(M))&=\diag(|M^T|^\alpha,|M|^\alpha),\\
 f_o(\mathcal D(M))&=
 \begin{pmatrix}0&\mathcal P_\alpha(M)\\
                 \mathcal P_\alpha(M)^T&0\end{pmatrix}.
\end{align*}
This already proves the rectangular error and cost, but generic real
node factoring would overcount the leading rank. The following exact
parity removes that loss.

Choose the OSE blockwise, so
$G_0=\diag(G_{0,M},I_m)$. Put
$P=\diag(I_N,-I_m)$ on physical space and
$Q_s=\diag(I_m,-I_m)$ on latent space. Then $PY=YQ_s$,
$J_sQ_s=-Q_sJ_s$, and $G_0,G-G_0$ commute with $Q_s$.
The symmetric finite polynomial of Section 3 consequently gives
\begin{equation}\label{eq:parityR}
 P\widehat R_{it}P=-\widehat R_{-it}
                       =-\overline{\widehat R_{it}}.
\end{equation}
Indeed $Q_sB_{0,z}Q_s=-B_{0,-z}$, so its degree-$j$ alternating word
acquires $(-1)^{j+1}$; the external resolvent scaling proves
\eqref{eq:parityR}.

The update $E_d=\mathcal D(\Delta)$ anticommutes with $P$. The computed
Woodbury node may be written
\[
 \widehat D=\widehat R E_d(I+\widehat R E_d)^{-1}\widehat R.
\]
This formula needs no inverse of $\widehat R$; the middle inverse
exists by the small-core determinant lemma. Substitution into
\eqref{eq:parityR} gives $P\widehat DP=-\overline{\widehat D}$.
Thus diagonal blocks are purely imaginary and off-diagonal blocks real.
Their complex ranks are at most $\rank E_d\le2r$. Dividing a purely
imaginary block by $i$, or taking a real off-diagonal block, preserves
that rank over the real field. Each extracted absolute or polar node
therefore has real rank at most $2r$, giving exactly
\eqref{eq:wholewidth}, not twice that coefficient.

A generic realification can first be formed, then thin QR and a small
exact rank factorization compress each block to $2r$ columns in
$O(Nr^2+r^3)$ work. Rectangular polar blocks use an ordinary two-factor
rank factorization, rather than a rectangular signed Cholesky.
Expanding a right absolute correction costs $O(m^2rL^2)$ and a polar
correction $O(NmrL^2)$, both already absorbed by the theorem. Left
absolute expansion has $N^2$ entries and its separately stated cost.
This completes Theorem~\ref{thm:rect}.

\section{A further PSD raw-base generality}

The positive-endpoint algorithm also applies to a promised PSD signed
raw base $A=XJX^T$ without materializing a new tall positive factor.
On the OSE event write $X=ZL$ with $Z=XL^\dagger$ and
$(1-\eta_0)I\preceq Z^TZ\preceq(1+\eta_0)I$.
Since $J$ restricted to $\operatorname{row}(X)$ is PSD,
$K_0=LJL^T\succeq0$. In fact this remains true on every sketch
outcome because $\operatorname{row}(L)\subseteq\operatorname{row}(X)$.
Compute small PSD Cholesky $K_0=HH^T$. The implicit positive factor
\[
 Y=XL^\dagger H=ZH,\qquad YY^T=A
\]
is exact on the OSE event. Its coarse Gram proxy is $H^TH$, and its
thin actions cost $O(Nmr+m^2r)$. Small support inverses and Cholesky
cost $O(m^3)$; no $N$-by-$m$ transformed factor is materialized.
Thus the same positive-base proof and operation count apply.

\section{Comparison, interpretation, and limits}

There is a direct same-guarantee explicit-Gram baseline. Form $X^TX$
in $O(Nm^2)$ conventional work, then use the same finite rational rule
with exact small shifted inverses and thin lifts. It has the same
intrinsic error, singular-endpoint scope, leading ranks, and conventional
cost $\widetilde O(Nm^2+m^3)$. The proved improvement here is
\[
 \widetilde O(Nm^2+m^3)\quad\longrightarrow\quad
 \widetilde O(Nmr+m^3).
\]
The ratio can approach $m/r$ when $Nr$ is sufficiently large compared
with $m^2$. For $N=\Theta(m)$ both cubic terms remain; for $r=m$ the
construction bounds also have the same order. These comparisons concern
these two proved realizations, not lower bounds for competing algorithms.

Root-exponential rank upper mechanisms already occur in
Birman--Koplienko--Solomyak~\cite{BKS}. Rational-function update algebra
is classical~\cite{BV}; rational Krylov update methods are developed
in Beckermann--Cortinovis--Kressner--Schweitzer~\cite{BCKS}. Their method
already handles general low-rank changes without taking the base
function as numerical input; its Hermitian positive-update square-root
approximant is PSD. Those properties are not unique to our construction.
Their printed bounds use rational approximation on enclosing spectral
sets, rather than the precise raw-factor/intrinsic-scale operation theorem
above. This is a scoped comparison of statements, not a claim that
earlier tools cannot be composed to imply our theorem.

Shmueli--Drineas--Avron~\cite{SDA} already treats root downdates and
Riccati formulations. Its printed guarantees assume SPD endpoints and
root/inverse-root access, with a stated iteration assumption in its cost.
Fasi--Higham--Liu~\cite{FHL} gives stronger exact rank-$r$ formulas for
scaled-identity bases; our uniform worst-case ranks do not supersede
that special structure. Sparse embeddings and preconditioned iterations
are established ingredients~\cite{NN}. There is substantial risk that
the raw operation theorem is a routine synthesis of known results.
A bounded primary-source review located no complete matching statement,
but neither novelty nor global state-of-the-art superiority is certified.

The sharp-rank claim uses the companion lower theorem, whose precise
matrix constants have their own classical scalar/Hankel antecedents
and priority qualifications~\cite{Sharp,Extensions}. It is an intrinsic
update-scale guarantee, not error relative to the possibly much smaller
correction norm. It gives fractional polar maps with $\alpha>0$,
not the discontinuous $\alpha=0$ polar factor. It does not give a
general quadratic matrix-product algorithm: squaring a nonzero threshold
$XJX^T-tI$ introduces the unavailable core $J(X^TX)J-2tJ$.

Exact arithmetic matters. The Gaussian polynomial primitive can have
severe coefficient cancellation, and the supplied factors may themselves
cancel severely. No numerical stability, finite-word complexity, SDP
solver runtime, or hardware speed conclusion is asserted.

\begin{thebibliography}{9}\raggedright
\bibitem{Sharp}
\emph{Exact small-exponent constants for fractional matrix maps}.
Working companion manuscript, November 2026.
\href{sharp-fractional-holder.pdf}{Local companion PDF}; source and
independent review record are in the same research workspace.

\bibitem{Extensions}
\emph{Factor four: finite-dimension crossover and structural significance}.
Working companion note, November 2026, Sections 13--14; independently
reviewed internally. Local file \texttt{factor-four-extensions.txt}.

\bibitem{NN}
J.~Nelson and H.~L.~Nguy\^en.
OSNAP: Faster numerical linear algebra algorithms via sparser subspace embeddings.
FOCS 2013; Theorem 9 of \url{https://arxiv.org/pdf/1211.1002}.

\bibitem{BKS}
M.~Sh.~Birman, L.~S.~Koplienko and M.~Z.~Solomyak.
Estimates of the spectrum of a difference of fractional powers of selfadjoint operators.
\emph{Izv. Vyssh. Uchebn. Zaved. Mat.} (1975), no.~3, 3--10.
\url{https://www.mathnet.ru/eng/ivm6330}.

\bibitem{BV}
D.~S.~Bernstein and C.~F.~Van Loan.
Rational matrix functions and rank-1 updates.
\emph{SIAM J. Matrix Anal. Appl.} \textbf{22} (2000), 145--154.
\href{https://doi.org/10.1137/S0895479898333636}{doi:10.1137/S0895479898333636}.

\bibitem{BCKS}
B.~Beckermann, A.~Cortinovis, D.~Kressner and M.~Schweitzer.
Low-rank updates of matrix functions II: Rational Krylov methods.
\emph{SIAM J. Numer. Anal.} \textbf{59} (2021), 1325--1347.
\url{https://arxiv.org/abs/2008.11501}.

\bibitem{SDA}
S.~Shmueli, P.~Drineas and H.~Avron.
Low-rank updates of matrix square roots.
\emph{Numer. Linear Algebra Appl.} \textbf{31} (2024), e2528;
first published online 2023.
\href{https://doi.org/10.1002/nla.2528}{doi:10.1002/nla.2528}.

\bibitem{FHL}
M.~Fasi, N.~J.~Higham and X.~Liu.
Computing the square root of a low-rank perturbation of the scaled identity matrix.
\emph{SIAM J. Matrix Anal. Appl.} \textbf{44} (2023), 156--174.
\href{https://doi.org/10.1137/22M1471559}{doi:10.1137/22M1471559}.
\end{thebibliography}
\end{document}
